\section{Methods: one certificate, end to end}\label{sec:methods}

%% Added 2026-09-14 under the 2026-09-13 ruling (decision sheet:
%% futon2 holes/labs/wm-contract/DECISION-SHEET-companion-reference-2026-09-13.md):
%% the companion promises resolve to futon-2026, and this paper gains a
%% Methods section carrying the Lean certificate once the relevant
%% certificate exists. It exists; this is it.

Where this paper's claims are validated, they are validated by
machine-checked certificate over pinned records, not by narrative.  This
section records one complete certificate chain --- for the machine's
pattern-retrieval act (\texttt{find}), the operation that selects catalog
patterns from the recorded library and must warrant each selection ---
so the reader can see concretely what is checked, over what record, and
on whose authority.  The companion paper owns the formal treatment and
the preregistered experimental report; what belongs here is the
certificate for the specimen this paper describes.\worknote{ADDED
2026-09-14: the certificate ruled for (2026-09-13) exists and is pinned
here --- the F1--F4 Lean closures plus the
\texttt{:wm/f2-reconciliation-certificate-v1} chain. Rows P001--P005
close.}

\paragraph{The record.}  The subject is a fixed, transcribed record of
retrieval behaviour: \texttt{futon3:checks/find-snatch.edn}, thirty-four
recorded rounds across six scenarios of \texttt{find} calls against an
eighteen-member recorded pattern repository, each selection carrying its
receipt, with one empty selection (scenario \texttt{g2/snatcher}, round
11) recorded as a typed absence.  The file's SHA-256 is
\texttt{c11673ea7164e90b10cc378ab6b2dfe14e545449d85e0dde70d5c2282e2430ce}.

\paragraph{The Lean leg.}  Four properties of that record are theorems in
\texttt{DarkTower/WarMachine/Holes.lean}: every selection lies inside the
recorded repository, with empty selection carried as typed absence
(\texttt{wmFindSnatchF1Containment}); every selected pattern carries a
receipt in the same recorded row (\texttt{wmFindSnatchF2Receipted});
every receipt is a structured antecedent with a warrant file rather than
a bare score (\texttt{wmFindSnatchF3NonSelfCertifying}); and each
scenario's declared zero-mass pattern is a repository member absent from
that scenario's selected union, so the selection claim is falsifiable in
principle and was not falsified (\texttt{wmFindSnatchF4Falsifiable}).
Each is proved by \texttt{decide} over the transcribed table --- the
kernel evaluates the predicate on all rows --- and none depends on
\texttt{sorry}.  The theorems closed under the J9 criterion on
2026-09-03; the file re-elaborated cleanly on 2026-09-14, its only two
\texttt{sorry} declarations being the deliberately open external
attestation holes H3 and H4 (theory-alignment claims about policy
precision that Lean cannot decide from a transcript), which do not touch
these four.  Each theorem has a rejecting negative control in
\texttt{futon3:checks/find\_snatch.clj} (\texttt{--negative-f1} through
\texttt{--negative-f4}): a corrupted counterpart record is refused, so
the checks are detectors, not constants.

\paragraph{The transcription leg.}  A \texttt{decide} proof is only as
good as the transcription it evaluates, so the chain's second leg binds
the transcribed table to the pinned file and the pinned file to history
that the transcriber does not control.  A committed expectation record
(\texttt{00-pin-expectation.edn}) declares the pin path and digest and
binds them to the futon3 repository's own git history --- commit
\texttt{e63eaef8}, 2026-09-08, authored five days before the
transcription and not by the transcriber, with the blob identity and a
re-derivation command any reader can run.  The verifier
(\texttt{f11\_f2\_reconcile.bb --check}) validates that expectation
fail-closed (exactly one EDN form, exact key set, declared path equal to
the path actually hashed, digest format, non-blank authority provenance
--- each with a refusing control), refuses any drift between the pinned
digest and the file on disk, then recomputes the full reconciliation
report and refuses field-level drift from the committed certificate
(\texttt{:wm/f2-reconciliation-certificate-v1}).  On 2026-09-14 the check
passes: the committed certificate describes exactly what recomputation
observes, and the records reconcile.

\paragraph{How the chain itself was built.}  The expectation-and-verifier
chain was authored inside the machine's own gated loop, not beside it.
An independent reviewer rejected the first version outright --- the
expectation was self-authored, so the digest check was reflexive --- and
returned a second round of findings against the repair (declared-path
and provenance validation gaps) before approving the amended chain; the
closing run then recorded a grounded change in a preregistered
machinery-test cohort with durable checkpoints (cohort 49, attempt 002,
2026-09-14).  Those runs are machinery evidence: they demonstrate the
loop that produced this certificate under review and refusal, and none
of them is, or is presented as, the preregistered qualifying run, which
remains outstanding (tracker row 28) and belongs to the companion
paper's experimental report.

\paragraph{Scope.}  This certificate establishes recorded-row invariants
of one pinned record and the integrity of the pin; it does not establish
runtime behaviour beyond that record, and on this record F3's
discrimination beyond F2 rests on its rejecting control, since every
recorded receipt is already a structured antecedent.  The open
theory-alignment holes (H3, H4) remain open by design and are reported,
not smoothed.
